\chapter{Installing}
\label{ch:install}

hdl-vet needs Python 3.11 or newer. It needs one other package, \code{pyslang}
version 12.1.0. The installer fetches it for you. The equivalence check needs one more
package. See the section about the extra \code{equiv} below.

You install hdl-vet from PyPI, from a checkout of its repository, or from a
wheel file. Use a virtual environment so that \code{pyslang} does not mix with
other packages.

\section{The systems it runs on}
\label{sec:platforms}

hdl-vet runs on Linux, on macOS and on Windows. CI tests it on all three: on Linux, on
macOS on Apple silicon, and on Windows.

macOS on Intel is not tested. It installs, because the built file that \code{pyslang}
gives for macOS serves both kinds of processor, and nothing in hdl-vet reads the kind.
Nobody runs the tests there, so this manual says what is true and no more.

No compiler is needed on any system. Both packages that hdl-vet needs come as built
files. \code{pyslang} has one for Linux, for macOS and for Windows. The Yosys of the
extra \code{equiv} is one file that holds WebAssembly. It runs anywhere, and it is the
same build on every system, so the equivalence check gives the same answer on all of
them.

On macOS, the Python that comes with the system is older than 3.11. Install a newer one,
for example with Homebrew or with uv, and make the virtual environment from that one.
\cmd{python3 --version} tells you which one you have.

\section{From PyPI}

\begin{enumerate}
\item Create and enter a virtual environment:
  \cmd{python -m venv .venv}, then activate it as your shell requires.
\item Install the package: \cmd{pip install hdl-vet}
\end{enumerate}

The name of the package is \code{hdl-vet}. The name of the command is \cmd{vet}.

\section{From a checkout}

\begin{enumerate}
\item Create and enter a virtual environment:
  \cmd{python -m venv .venv}, then activate it as your shell requires.
\item Go to the folder of the checkout, the one that holds \file{pyproject.toml}.
\item Install it: \cmd{pip install .}
\end{enumerate}

For work on hdl-vet itself, \cmd{pip install -e .} makes an editable install.

\section{From a wheel}

A wheel is one file with the extension \file{.whl}. If you have one,
install it with \cmd{pip install FILE.whl}. To make a wheel from a checkout,
use any build tool that reads \file{pyproject.toml}, for example
\cmd{uv build --wheel}. The wheel appears in the folder \file{dist}.

\section{The extra for the equivalence check}
\label{sec:extra}

\apiname{Module.override_net} changes logic. The sign-off proves the change with Yosys
(section~\ref{sec:equiv}). Yosys comes from the package \code{yowasp-yosys}. hdl-vet pins
one version of it, because a new version of Yosys can change an answer.

The extra \code{equiv} installs this package. Give the name of the extra in square
brackets. Quote the argument, so that your shell does not read the brackets:

\begin{itemize}
\item \cmd{pip install "hdl-vet[equiv]"} installs the release from PyPI.
\item \cmd{pip install ".[equiv]"} installs a checkout.
\end{itemize}

Without the extra, hdl-vet works as before. The one exception is a logic operation. Without
Yosys, \apiname{Module.override_net} refuses the call unless it has a waiver. The message
names \code{hdl-vet[equiv]}.

\section{Check the installation}

Both of these print the program name and the version, for example
\code{hdl-vet 1.2.3}, and nothing else:

\begin{itemize}
\item \cmd{vet --version}
\item \cmd{python -m hdl_vet --version}
\end{itemize}

The commands \cmd{vet} and \cmd{python -m hdl_vet} are the same program. The
transcripts in this manual use \cmd{vet}.
